\documentclass[addpoints,10pt]{exam}
\usepackage[utf8]{inputenc}
\usepackage[T1]{fontenc}
\usepackage[spanish]{babel}
\usepackage{times}
\usepackage{enumerate}
\usepackage{graphicx}
\usepackage{algorithm2e}
\usepackage{stfloats}
\usepackage{multirow}
\usepackage{array}
\usepackage{vhistory}
\usepackage{amsmath}
\usepackage{tcolorbox}
\usepackage{fancyvrb}
\usepackage[lighttt]{lmodern}
\usepackage{xcolor}
\usepackage{listings}
\lstset{
	basicstyle=\small\ttfamily,
	stringstyle=\ttfamily,
	commentstyle=\color{gray},
	showstringspaces=false,
	keywordstyle=\bfseries
}
\usepackage{multicol,caption,capt-of}
\usepackage[left=1.5cm,right=1.5cm,top=1.5cm,bottom=1.5cm]{geometry}

\begin{document}
	
	\pagestyle{headandfoot}
	\runningheadrule
	\firstpageheader{Programación}{Evaluación Unidad 4 - Estructuras de datos}{2026}
	\runningheader{Programación}{Evaluación Unidad 4, Page \thepage\ of \numpages}{2026}
	\firstpagefooter{Escuela de Ingenierías}{UPB}{Ingeniería Aeronáutica}
	\runningfooter{Escuela de Ingenierías}{UPB}{Ingeniería Aeronáutica}
	
	\begin{center}
		\begin{tabular}{|m{4cm}|m{8cm}|m{6cm}|}
			\hline
			\multirow{5}{*}{\includegraphics[width=1\linewidth]{Logo-UPB-2022.png}}
			&\multicolumn{2}{c|}{\textbf{Universidad Pontificia Bolivariana - Sede Medellín}}\\ \cline{2-3}
			&Curso: \textbf{Programación}&Profesor: \textbf{Henry Andrade, Ph.D.}\\ \cline{2-3}
			&Duración: \textbf{100 minutos}&Total: \textbf{100 puntos}\\
			&\textbf{Nombre:}\ &\textbf{ID:} \\\cline{2-3} \hline
		\end{tabular}
	\end{center}
	
	\tcbset{colback=red!5!white, colframe=red!75!black, fonttitle=\bfseries}
	\begin{tcolorbox}[title=Nota Importante]
		Lea cuidadosamente cada pregunta. Responda utilizando únicamente conceptos y estructuras trabajados en clase. No se permite el uso de \emph{list comprehensions}, librerías externas ni acceso a Internet durante la prueba. En las preguntas de código, respete la sintaxis e indentación de Python.
	\end{tcolorbox}
	
	\begin{multicols}{2}
		\begin{questions}
			\boxedpoints
			\pointname{ Puntos}
			\qformat{Pregunta \thequestion \dotfill \thepoints}
			
			\section*{1. Fundamentos y conceptos (30 puntos)}
			
			\question[10]
			\textbf{Emparejamiento.} Escriba en la columna central la letra de la descripción que corresponde a cada concepto.
			
			\noindent
			\begin{center}
				\begin{tabular}{|l|c|p{4.6cm}|}
					\hline
					\textbf{Concepto} & \textbf{Letra} & \textbf{Descripción}\\ \hline
					1. Lista & & A. Propiedad que indica si un objeto puede modificarse después de creado.\\ \hline
					2. Tupla & & B. Colección ordenada e inmutable; admite desempaquetado.\\ \hline
					3. Diccionario & & C. Función que retorna el identificador único de un objeto.\\ \hline
					4. Mutabilidad & & D. Colección de pares clave-valor con claves únicas.\\ \hline
					5. \texttt{id()} & & E. Colección ordenada y mutable que admite elementos repetidos.\\ \hline
				\end{tabular}
			\end{center}
			
			\question[10]
			\textbf{Completar frases.} Complete cada enunciado con una palabra del banco de palabras. Cada palabra se usa máximo una vez.
			
			\begin{center}
				\fbox{\texttt{mutable \ inmutable \ etiqueta \ key \ value \ zip \ enumerate}}
			\end{center}
			
			\begin{enumerate}
				\item En Python, una variable funciona como una \rule{2cm}{0.4pt} que apunta a un objeto.
				\item Un \texttt{string} es un objeto \rule{2cm}{0.4pt}: al "modificarlo" se crea uno nuevo.
				\item En un diccionario, el elemento que identifica de forma única a un dato es la \rule{2cm}{0.4pt}.
				\item Para recorrer dos listas en paralelo se utiliza la función \rule{2cm}{0.4pt}.
				\item Para obtener el índice y el valor al recorrer una lista se utiliza \rule{2cm}{0.4pt}.
			\end{enumerate}
			
			\question[10]
			\textbf{Etiquetas y mutabilidad.} Analice los siguientes fragmentos de código.
			
\begin{lstlisting}
altitud = 10000
elevacion = altitud
elevacion = 9500
print(altitud, elevacion)
\end{lstlisting}

\begin{lstlisting}
def agregar(tanques, litros):
   tanques.append(litros)

combustible = [1000, 1200]
agregar(combustible, 500)
print(combustible)
\end{lstlisting}
			
			\begin{enumerate}
				\item Escriba la salida de cada fragmento.
				\item Explique por qué \texttt{altitud} no cambió, pero \texttt{combustible} sí fue modificada dentro de la función.
			\end{enumerate}
			
			\makeemptybox{4cm}
			
			\section*{2. Listas, tuplas y diccionarios (45 puntos)}
			
			\question[10]
			\textbf{Listas: métodos e índices.} Complete el código con el método o índice que corresponda según el comentario.
			
\begin{lstlisting}
componentes = ["alas", "fuselaje", "motores"]

componentes.________("tren de aterrizaje")  # agrega al final
print(componentes[____])    # imprime el ultimo elemento
componentes.________("alas")  # elimina "alas"
pos = componentes.________("motores")  # posicion de "motores"
\end{lstlisting}
			
			\question[10]
			\textbf{Tuplas.} Complete el código teniendo en cuenta que las tuplas son inmutables.
			
\begin{lstlisting}
coordenada = (33.9425, -118.4081)
lat, lon = ____________  # desempaquetado

avion = ("Boeing 787", "Dreamliner", 2009, 242)
print(avion[____])       # imprime 242

rumbo = (270____)        # tupla de un elemento
\end{lstlisting}
			
			\begin{enumerate}
				\item Complete los espacios en blanco.
				\item ¿Qué ocurre si se ejecuta \texttt{avion[0] = "Airbus"}? Justifique.
			\end{enumerate}
			
			\makeemptybox{2cm}
			
			\question[15]
			\textbf{Estructuras anidadas (diccionarios y listas).} Una aerolínea almacena su flota como un diccionario cuya clave es la matrícula y cuyo valor es otro diccionario con el modelo y una lista de componentes. Complete las líneas faltantes para imprimir los componentes que superaron su límite de horas.
			
\begin{lstlisting}
flota = {
  "HK-4532": {"modelo": "A320", "componentes": [
     {"nombre": "Motor izq", "horas": 1200, "limite": 1000},
     {"nombre": "Tren", "horas": 900, "limite": 4500}]},
  "HK-1020": {"modelo": "ATR72", "componentes": [
     {"nombre": "Alabes", "horas": 3100, "limite": 3000}]}
}

for matricula in flota:
   for comp in ____________________________:
      if ____________________________:
         print(matricula, comp["nombre"],
               "requiere mantenimiento")
\end{lstlisting}
			
			\textbf{Escriba además} la salida que produce el programa completo:
			\makeemptybox{2cm}
			
			\question[10]
			\textbf{Deducción de funcionamiento.} Analice el siguiente fragmento y escriba lo que imprime cada \texttt{print}.
			
\begin{lstlisting}
meteo = {
  "KLAX": {"temp": 22, "viento": 10},
  "KORD": {"temp": 5, "viento": 25}
}
meteo["KORD"]["temp"] = 3
meteo["KMIA"] = {"temp": 28, "viento": 8}

print(len(meteo))
print(meteo.get("KJFK", "sin datos"))
print(meteo["KORD"]["temp"])
\end{lstlisting}
			
			Salida:
			\makeemptybox{2.5cm}
			
			\section*{3. Auditoría de código y reto (25 puntos)}
			
			\question[15]
			\textbf{Auditoría de código (reto de mantenimiento).} Un grupo desarrolló el siguiente programa que registra 3 aeronaves y reporta las que superan las 5000 horas de vuelo.
			
\begin{lstlisting}
flota = {}
for i in range(3):
   matricula = input("Matricula: ")
   horas = input("Horas de vuelo: ")
   flota[matricula] = [horas]
   flota[matricula].append(input("Modelo: "))

for matricula in flota:
   if flota[matricula][0] > 5000:
      print(matricula, "requiere mantenimiento")
\end{lstlisting}
			
			\begin{enumerate}
				\item Identifique el error que impedirá que el programa funcione correctamente y explique por qué ocurre. (5 puntos)
				\item Escriba la línea corregida. (4 puntos)
				\item Proponga una mejora en la estructura de datos para que el código sea más legible y justifíquela. (6 puntos)
			\end{enumerate}
			
			\makeemptybox{5cm}
			
			\question[10]
			\textbf{Auditoría: copias de listas.} Un estudiante escribió el siguiente código para conservar una copia de seguridad de las lecturas originales, y afirma que \texttt{respaldo} no se verá afectada por los cambios.
			
\begin{lstlisting}
lecturas = [1000, 1200, 800]
respaldo = lecturas
lecturas.append(500)
print(respaldo)
\end{lstlisting}
			
			\begin{enumerate}
				\item Indique qué imprime realmente el programa y explique por qué la afirmación del estudiante es incorrecta. (5 puntos)
				\item Escriba la línea corregida para que \texttt{respaldo} sea independiente. (5 puntos)
			\end{enumerate}
			
			\makeemptybox{3.5cm}
		
		\end{questions}
	\end{multicols}
	
%	\begin{center}
%		\gradetable[h][questions]
%	\end{center}
	
\end{document}
